\section{Cronograma da pesquisa}

A seguir apresenta-se a estimativa preliminar das atividades de pesquisa e de seus respectivos prazos. O cronograma toma maio de 2026 como mês inicial e deverá ser ajustado conforme o calendário oficial do programa, a definição de orientação, a qualificação e a disponibilidade de avaliadores humanos.

\begin{landscape}
\begingroup
\scriptsize
\setlength{\tabcolsep}{2.5pt}
\newcommand{\mes}[1]{\rotatebox{90}{\textbf{#1}}}
\begin{longtable}{@{}p{0.34\linewidth}*{10}{>{\centering\arraybackslash}p{0.045\linewidth}}@{}}
\caption{Descritivo preliminar de atividades e prazos de entrega}
\label{tab:cronograma}\\
\toprule
\textbf{Atividades de pesquisa} & \mes{mai/26} & \mes{jun/26} & \mes{jul/26} & \mes{ago/26} & \mes{set/26} & \mes{out/26} & \mes{nov/26} & \mes{dez/26} & \mes{jan/27} & \mes{fev/27} \\
\midrule
\endfirsthead
\toprule
\textbf{Atividades de pesquisa} & \mes{mai/26} & \mes{jun/26} & \mes{jul/26} & \mes{ago/26} & \mes{set/26} & \mes{out/26} & \mes{nov/26} & \mes{dez/26} & \mes{jan/27} & \mes{fev/27} \\
\midrule
\endhead
Início da orientação e refinamento do problema & X & X &  &  &  &  &  &  &  &  \\
Revisão de literatura e matriz teórica & X & X & X &  &  &  &  &  &  &  \\
Definição do protocolo de avaliação &  & X & X &  &  &  &  &  &  &  \\
Coleta, tratamento e organização do corpus &  & X & X & X &  &  &  &  &  &  \\
Construção da bateria de perguntas e conjunto de referência &  &  & X & X &  &  &  &  &  &  \\
Piloto metodológico e ajustes do protocolo &  &  &  & X & X &  &  &  &  &  \\
Exame de qualificação &  &  &  & X &  &  &  &  &  &  \\
Aplicação do protocolo avaliativo ao sistema RAG &  &  &  &  & X & X &  &  &  &  \\
Avaliação humana e avaliação baseada no paradigma \textit{LLM-as-a-judge} &  &  &  &  &  & X & X &  &  &  \\
Análise dos dados e interpretação dos resultados &  &  &  &  &  &  & X & X &  &  \\
Síntese dos resultados, discussão e recomendações &  &  &  &  &  &  &  & X & X &  \\
Redação e revisão dos capítulos teóricos e metodológicos & X & X & X & X &  &  &  &  &  &  \\
Redação dos resultados, discussão e conclusão &  &  &  &  &  &  & X & X & X &  \\
Depósito da versão final da dissertação &  &  &  &  &  &  &  &  & X &  \\
Previsão de defesa da dissertação &  &  &  &  &  &  &  &  &  & X \\
Elaboração de artigo científico derivado &  &  &  &  &  &  &  &  & X & X \\
\bottomrule
\end{longtable}
\endgroup

\noindent Fonte: elaborado pelo autor (2026).
\end{landscape}
